
\(\Projz\) sur la catégorie des \(\underline{B}\)-Modules gradués est compatible au produit tensoriel; il en est encore de même du foncteur obtenu sur les catégories dérivées.


2.4. Pour tout entier $k$, et tout $\underline{O}_X$-Module $\underline{L}$, posons :
\begin{equation}\tag{2.4.1}
  F_k(L)=\bigoplus_{n\geq k} f_*\bigl(\underline{L}(n)\bigr).
\end{equation}
Le foncteur $F_k:\mathcal A\to\mathcal C$ ainsi défini est exact à gauche. Comme la catégorie des $\underline{O}_X$-Modules a assez d'injectifs, on peut définir un foncteur dérivé :
\begin{equation}\tag{2.4.2}
  \underline{\underline{\mathrm R}} F_k:D^+(\mathcal A)\longrightarrow D^+(\mathcal C).
\end{equation}

Si $\underline{L}$ est un $\underline{O}_X$-Module localement libre de type fini, il existe un entier $k_0$ tel que pour $k\geq k_0$, $\underline{L}$ soit $F_k$-acyclique. Il suffit pour le voir de montrer que, pour $n$ assez grand, $\underline{L}$ est acyclique pour les foncteurs $L\leadsto f_*\bigl(\underline{L}(n)\bigr)$. Comme la tensorisation par $\underline{O}_X(n)$ est un foncteur exact, cela résulte immédiatement de la suite spectrale des foncteurs composés et de 1.3.

2.5. On se propose de construire pour tout $\underline{O}_X$-Module $\underline{F}$ un homomorphisme fonctoriel en $\underline{F}$ :
\begin{equation}\tag{2.5.1}
  \beta:\Projz\bigl(F_k(\underline{F})\bigr)\longrightarrow \underline{F}.
\end{equation}
Pour cela on définit la restriction de $\beta$ aux $D_+(U,s)$ (2.3.2). Or par construction on a, en posant $\underline{M}=F_k(\underline{F})$ :
\[
  \Projz(\underline{M})|D_+(U,s)
  =\bar f^{*}_{U,s}\bigl(\underline{M}_{(s),U}\bigr).
\]
Pour définir un homomorphisme~:
\[
  \bar f^{*}_{U,s}\bigl(\underline{M}_{(s),U}\bigr)\longrightarrow \underline{F}|D_+(U,s),
\]
